\documentclass[11pt]{article}

\usepackage[margin=0.82in,headheight=22pt]{geometry}
\usepackage{amsmath,amssymb}
\usepackage{fontspec}
\setmainfont{TeX Gyre Pagella}
\setsansfont{TeX Gyre Heros}
\setmonofont{Courier New}[Scale=0.84]
\usepackage{microtype}
\usepackage{xcolor}
\usepackage{booktabs,longtable,array,colortbl}
\usepackage{enumitem}
\usepackage{listings}
\usepackage{fancyhdr}
\usepackage{titlesec}
\usepackage{hyperref}
\usepackage{bookmark}

\definecolor{navy}{RGB}{24,55,95}
\definecolor{teal}{RGB}{20,112,104}
\definecolor{codebg}{RGB}{246,248,250}
\definecolor{rulegray}{RGB}{210,216,222}
\definecolor{codecomment}{RGB}{80,100,90}

\hypersetup{
  colorlinks=true,
  linkcolor=navy,
  urlcolor=teal,
  pdftitle={STAT452 Practical Midterm: Complete R Solution},
  pdfauthor={STAT452 Study Guide}
}

\pagestyle{fancy}
\fancyhf{}
\lhead{\sffamily\small STAT452 Practical Midterm}
\rhead{\sffamily\small Complete R Solution}
\cfoot{\sffamily\thepage}
\renewcommand{\headrulewidth}{0.4pt}
\renewcommand{\headrule}{\hbox to\headwidth{\color{navy!45}\leaders\hrule height \headrulewidth\hfill}}

\titleformat{\section}{\sffamily\LARGE\bfseries\color{navy}}{\thesection}{0.65em}{}
\titleformat{\subsection}{\sffamily\Large\bfseries\color{teal}}{\thesubsection}{0.65em}{}
\titleformat{\subsubsection}{\sffamily\large\bfseries\color{navy!85}}{\thesubsubsection}{0.65em}{}
\titlespacing*{\section}{0pt}{2em}{0.75em}
\titlespacing*{\subsection}{0pt}{1.45em}{0.55em}
\titlespacing*{\subsubsection}{0pt}{1.15em}{0.4em}

\setlength{\parindent}{0pt}
\setlength{\parskip}{0.5em}
\setlength{\emergencystretch}{1.5em}
\setlist{leftmargin=1.55em,itemsep=0.25em,topsep=0.35em}
\renewcommand{\arraystretch}{1.18}
\setlength{\LTpre}{0.4em}
\setlength{\LTpost}{0.6em}

\lstset{
  language=R,
  basicstyle=\ttfamily\small,
  keywordstyle=\color{navy}\bfseries,
  commentstyle=\color{codecomment}\itshape,
  stringstyle=\color{teal!80!black},
  backgroundcolor=\color{codebg},
  frame=single,
  rulecolor=\color{rulegray},
  breaklines=true,
  columns=fullflexible,
  keepspaces=true,
  showstringspaces=false,
  xleftmargin=0.45em,
  xrightmargin=0.45em,
  aboveskip=0.6em,
  belowskip=0.7em
}

\begin{document}
\hypersetup{pageanchor=false}

\begin{titlepage}
  \centering
  \vspace*{1.0in}
  {\sffamily\color{navy}\Huge\bfseries STAT452 Practical Midterm\par}
  \vspace{0.25in}
  {\sffamily\color{teal}\LARGE Complete Solution Using R\par}
  \vspace{0.5in}
  {\color{navy!45}\rule{0.78\textwidth}{1.1pt}\par}
  \vspace{0.5in}
  {\Large OLS, Diagnostics, Ridge, LASSO, and Prediction\par}
  \vfill
  \begin{minipage}{0.82\textwidth}
    \colorbox{navy!7}{\parbox{0.94\linewidth}{
      \textbf{R-only workflow.} Every requested calculation is performed with
      R. The document shows executable code, the important output, and the
      conclusion that should be written for each exam part. No calculator or
      hand-computation route is used.
    }}
  \end{minipage}
  \vfill
  {\sffamily\small Based on \texttt{Midterm\_Practical.pdf},
  \texttt{seatpos.csv}, \texttt{newdata.csv}, and Notes Chapters 1 and 3\par}
\end{titlepage}

\hypersetup{pageanchor=true}
\pagenumbering{roman}
\tableofcontents
\clearpage
\pagenumbering{arabic}

\section{Reproducible setup}

The practical dataset contains 38 drivers, eight predictors, and the response
\texttt{hipcenter}. Chapter 1 of the Notes supplies the OLS and residual-analysis
workflow; Chapter 3 supplies cross-validation, ridge, and LASSO.

Run the companion script from the repository root:

\begin{lstlisting}
source("R_training/midterm/practical_test_solution.R")
\end{lstlisting}

Only \texttt{glmnet} is required beyond base R:

\begin{lstlisting}
install.packages("glmnet")  # Run once only, if needed
library(glmnet)
\end{lstlisting}

\subsection{Read and validate the files}

\begin{lstlisting}
seatpos <- read.csv("R_training/midterm/seatpos.csv")
newdata <- read.csv("R_training/midterm/newdata.csv")

dim(seatpos)                 # 38 9
names(seatpos)
str(seatpos)
summary(seatpos)
colSums(is.na(seatpos))      # all zero

stopifnot(nrow(seatpos) == 38, ncol(seatpos) == 9)
stopifnot(identical(
  names(newdata),
  setdiff(names(seatpos), "hipcenter")
))
\end{lstlisting}

These checks prevent silent errors from a wrong file, altered column name, or
missing data. The training data contain no missing values.

\section{Part (a): Full OLS model and residual analysis}

\subsection{Fit the full model}

\begin{lstlisting}
full_ols <- lm(hipcenter ~ ., data = seatpos)
summary(full_ols)
confint(full_ols)
\end{lstlisting}

The coefficient output is:

\begin{center}
\begin{tabular}{lrrr}
\toprule
\rowcolor{navy!10}\textbf{Term}&\textbf{Estimate}&\textbf{Std. error}&\textbf{$p$-value}\\
\midrule
Intercept&436.4321&166.5716&0.014\\
Age&0.7757&0.5703&0.184\\
Weight&0.0263&0.3310&0.937\\
HtShoes&-2.6924&9.7530&0.784\\
Ht&0.6013&10.1299&0.953\\
Seated&0.5338&3.7619&0.888\\
Arm&-1.3281&3.9002&0.736\\
Thigh&-1.1431&2.6600&0.671\\
Leg&-6.4390&4.7139&0.182\\
\bottomrule
\end{tabular}
\end{center}

The fitted R model is
\begin{align*}
\widehat{\texttt{hipcenter}}
={}&436.4321+0.7757\,\texttt{Age}
+0.0263\,\texttt{Weight}\\
&-2.6924\,\texttt{HtShoes}
+0.6013\,\texttt{Ht}
+0.5338\,\texttt{Seated}\\
&-1.3281\,\texttt{Arm}
-1.1431\,\texttt{Thigh}
-6.4390\,\texttt{Leg}.
\end{align*}

Each slope is interpreted while the other seven predictors are held fixed. For
example, the fitted coordinate increases by $0.776$ mm per additional year of
age, holding weight and all body dimensions fixed. A one-centimeter increase in
lower-leg length is associated with a $6.439$ mm decrease in the fitted
coordinate, holding the other predictors fixed.

The intercept is the fitted coordinate when every predictor equals zero, which
is far outside the observed anthropometric range and has no useful physical
interpretation.

\subsection{Model fit and the apparent coefficient paradox}

R reports
\[
R^2=0.687,\qquad R^2_{\mathrm{adj}}=0.600,
\qquad F(8,29)=7.94,\qquad p=1.31\times10^{-5}.
\]
The predictors are highly useful jointly, yet none of the eight individual
slopes has $p<0.05$. This is expected under severe multicollinearity: correlated
predictors inflate individual standard errors even when the group explains the
response well.

\subsection{Standard residual plots in R}

\begin{lstlisting}
old_par <- par(mfrow = c(2, 2))
plot(full_ols)
par(old_par)
\end{lstlisting}

Interpret the four R plots as follows:

\begin{itemize}
  \item \textbf{Residuals vs Fitted}: check linearity and look for a random cloud
  around zero.
  \item \textbf{Normal Q-Q}: points near the reference line support normality.
  \item \textbf{Scale-Location}: a roughly horizontal band supports constant
  error variance.
  \item \textbf{Residuals vs Leverage}: identify high-leverage observations with
  large Cook's distance.
\end{itemize}

\subsection{Normality and constant variance in R}

\begin{lstlisting}
shapiro.test(resid(full_ols))

# Base-R Breusch-Pagan LM test
predictor_data <- seatpos[setdiff(names(seatpos), "hipcenter")]
bp_aux <- lm(I(resid(full_ols)^2) ~ ., data = predictor_data)
bp_stat <- nrow(seatpos) * summary(bp_aux)$r.squared
bp_df <- length(coef(full_ols)) - 1
bp_p <- pchisq(bp_stat, df = bp_df, lower.tail = FALSE)
c(BP = bp_stat, df = bp_df, p_value = bp_p)
\end{lstlisting}

The Shapiro-Wilk result is $W=0.9715$, $p=0.434$. R therefore gives no evidence
against residual normality at $\alpha=0.05$. The Breusch-Pagan calculation gives
$BP=14.037$, $df=8$, $p=0.0808$, so R does not detect significant
heteroscedasticity at $5\%$, although the scale-location plot should still be
checked.

\subsection{Influence diagnostics in R}

\begin{lstlisting}
influence_table <- data.frame(
  row = seq_len(nrow(seatpos)),
  studentized_residual = rstudent(full_ols),
  leverage = hatvalues(full_ols),
  cooks_distance = cooks.distance(full_ols)
)

influence_table <- influence_table[
  order(influence_table$cooks_distance, decreasing = TRUE),
]
head(influence_table, 5)

4 / nrow(seatpos)                         # Cook's D guide
2 * length(coef(full_ols)) / nrow(seatpos) # leverage guide
\end{lstlisting}

\begin{center}
\begin{tabular}{rrrr}
\toprule
\rowcolor{navy!10}\textbf{Row}&\textbf{Studentized residual}&\textbf{Leverage}&\textbf{Cook's $D$}\\
\midrule
31& 2.390&0.560&0.695\\
23& 1.938&0.270&0.141\\
35&-2.323&0.183&0.117\\
33& 1.528&0.248&0.082\\
17&-0.914&0.401&0.063\\
\bottomrule
\end{tabular}
\end{center}

Row 31 exceeds the rough Cook's-distance guide $4/n=0.105$ and leverage guide
$2k/n=0.474$, where $k=9$ fitted coefficients including the intercept. It should
be checked and used in a sensitivity analysis, not automatically deleted.

\subsection{Multicollinearity in R}

\begin{lstlisting}
predictors <- setdiff(names(seatpos), "hipcenter")

vif <- vapply(predictors, function(variable) {
  others <- setdiff(predictors, variable)
  auxiliary <- lm(
    reformulate(others, response = variable),
    data = seatpos
  )
  1 / (1 - summary(auxiliary)$r.squared)
}, numeric(1))

sort(vif, decreasing = TRUE)
kappa(scale(seatpos[predictors]), exact = TRUE)
\end{lstlisting}

\begin{center}
\begin{tabular}{lr@{\qquad}lr}
\toprule
\rowcolor{navy!10}\textbf{Variable}&\textbf{VIF}&\textbf{Variable}&\textbf{VIF}\\
\midrule
Ht&333.138&Arm&4.496\\
HtShoes&307.429&Weight&3.647\\
Seated&8.951&Thigh&2.763\\
Leg&6.694&Age&1.998\\
\bottomrule
\end{tabular}
\end{center}

The standardized-predictor condition number is $59.77$. The extreme VIF values
for \texttt{Ht} and \texttt{HtShoes} confirm severe collinearity, explaining the
unstable OLS slopes and motivating ridge/LASSO.

\section{Part (b): Full model versus noise-only model}

Fit the intercept-only model and ask R for the nested-model partial $F$-test:

\begin{lstlisting}
null_ols <- lm(hipcenter ~ 1, data = seatpos)
anova(null_ols, full_ols)
\end{lstlisting}

\begin{center}
\begin{tabular}{lrrrrr}
\toprule
\rowcolor{navy!10}\textbf{Model}&\textbf{Res.Df}&\textbf{RSS}&\textbf{Df}&\textbf{$F$}&\textbf{$p$}\\
\midrule
Intercept only&37&131639&&&\\
Full model&29&41262&8&7.94&$1.31\times10^{-5}$\\
\bottomrule
\end{tabular}
\end{center}

At $\alpha=0.05$, R gives $p<0.05$, so reject the null hypothesis that all
eight slopes equal zero. The full model improves significantly on the
noise-only model.

\section{Parts (c)--(f): Shared \texttt{glmnet} setup}

\texttt{glmnet} requires a numeric predictor matrix rather than a data frame.
The same fold IDs must be used for ridge and LASSO to make the comparison fair.

\begin{lstlisting}
library(glmnet)

X <- model.matrix(hipcenter ~ ., data = seatpos)[, -1, drop = FALSE]
y <- seatpos$hipcenter
X_new <- model.matrix(~ ., data = newdata)[, -1, drop = FALSE]

lambda_grid <- seq(1e-3, 50, by = 1e-3)
set.seed(452)
fold_id <- sample(rep(1:10, length.out = nrow(X)))
\end{lstlisting}

The explicit \texttt{fold\_id} makes the answer reproducible. A different random
fold allocation can produce a slightly different optimal penalty.

\section{Part (c): Full ridge regression}

Set \texttt{alpha = 0}. The exam's 50,000-point grid is unusually dense, so the
current \texttt{glmnet} version may warn near the smallest penalties. The selected
penalty occurs safely before that path limit.

\begin{lstlisting}
ridge_cv <- suppressWarnings(cv.glmnet(
  x = X,
  y = y,
  alpha = 0,
  lambda = lambda_grid,
  foldid = fold_id,
  type.measure = "mse",
  standardize = TRUE
))

ridge_cv$lambda.min
ridge_cv$lambda.1se
coef(ridge_cv, s = "lambda.min")
plot(ridge_cv)
abline(v = log(ridge_cv$lambda.min), col = "red", lty = 2)
\end{lstlisting}

R returns
\[
\lambda_{\min}=36.521,
\qquad
\lambda_{1\mathrm{se}}=50.000.
\]
Because $\lambda_{1\mathrm{se}}$ equals the grid's upper boundary, a real
analysis using the 1-SE rule should expand the upper end beyond 50.

At $\lambda_{\min}$, R reports:

\begin{center}
\begin{tabular}{lr@{\qquad}lr}
\toprule
\rowcolor{navy!10}\textbf{Term}&\textbf{Ridge}&\textbf{Term}&\textbf{Ridge}\\
\midrule
Intercept&399.992503&Seated&-1.203317\\
Age&0.473581&Arm&-1.354624\\
Weight&-0.106186&Thigh&-1.304223\\
HtShoes&-0.738541&Leg&-3.109081\\
Ht&-0.737116&&\\
\bottomrule
\end{tabular}
\end{center}

Ridge keeps every slope nonzero but shrinks the unstable OLS estimates. This
matches Chapter 3: the $L_2$ penalty trades a small amount of bias for lower
variance under multicollinearity. \texttt{glmnet} standardizes predictors
internally, leaves the intercept unpenalised, and returns coefficients on the
original variable scales.

\section{Part (d): Ridge predictions}

\begin{lstlisting}
ridge_prediction <- drop(predict(
  ridge_cv,
  newx = X_new,
  s = "lambda.min"
))
ridge_prediction
\end{lstlisting}

\begin{center}
\begin{tabular}{rr}
\toprule
\rowcolor{navy!10}\textbf{New row}&\textbf{Ridge prediction (mm)}\\
\midrule
1&-139.711\\
2&-112.954\\
3&-219.990\\
4&-149.057\\
5&-228.157\\
6&-102.596\\
\bottomrule
\end{tabular}
\end{center}

\section{Part (e): Full LASSO model}

Set \texttt{alpha = 1} and reuse the same grid and folds:

\begin{lstlisting}
lasso_cv <- cv.glmnet(
  x = X,
  y = y,
  alpha = 1,
  lambda = lambda_grid,
  foldid = fold_id,
  type.measure = "mse",
  standardize = TRUE
)

lasso_cv$lambda.min
lasso_cv$lambda.1se
lasso_coef <- as.matrix(coef(lasso_cv, s = "lambda.min"))
lasso_coef

nonzero <- rownames(lasso_coef)[
  lasso_coef[, 1] != 0 & rownames(lasso_coef) != "(Intercept)"
]
nonzero
\end{lstlisting}

R returns
\[
\lambda_{\min}=5.748,
\qquad
\lambda_{1\mathrm{se}}=20.288.
\]
At $\lambda_{\min}$, the nonzero predictors are
\[
\boxed{\texttt{Age},\quad\texttt{Ht},\quad\texttt{Leg}.}
\]

\begin{center}
\begin{tabular}{lrrr}
\toprule
\rowcolor{navy!10}\textbf{Term}&\textbf{OLS}&\textbf{Ridge}&\textbf{LASSO}\\
\midrule
Intercept&436.432128&399.992503&397.696648\\
Age&0.775716&0.473581&0.221804\\
Weight&0.026313&-0.106186&0\\
HtShoes&-2.692408&-0.738541&0\\
Ht&0.601345&-0.737116&-2.200605\\
Seated&0.533752&-1.203317&0\\
Arm&-1.328069&-1.354624&0\\
Thigh&-1.143119&-1.304223&0\\
Leg&-6.439046&-3.109081&-5.468787\\
\bottomrule
\end{tabular}
\end{center}

The LASSO solution is sparse: five slopes are exactly zero. With strongly
correlated height measurements, LASSO may select one representative and discard
another, so the selected set can vary with the fold split. Ridge is usually more
stable when all correlated measurements carry signal; LASSO is more useful when
a compact predictive model is required.

For an optional post-LASSO comparison in R:

\begin{lstlisting}
post_lasso_ols <- lm(
  reformulate(nonzero, response = "hipcenter"),
  data = seatpos
)
summary(post_lasso_ols)
\end{lstlisting}

The refitted OLS slopes are approximately $0.581$ for Age, $-2.325$ for Ht, and
$-6.739$ for Leg. These are unpenalised estimates and therefore have larger
magnitudes than the LASSO coefficients. The refit is a different estimator; it
must not be labelled as the LASSO model.

\section{Part (f): Compare ridge and LASSO predictions}

\begin{lstlisting}
lasso_prediction <- drop(predict(
  lasso_cv,
  newx = X_new,
  s = "lambda.min"
))

prediction_comparison <- data.frame(
  row = seq_len(nrow(newdata)),
  ridge = ridge_prediction,
  LASSO = lasso_prediction,
  difference = ridge_prediction - lasso_prediction
)
prediction_comparison
\end{lstlisting}

\begin{center}
\begin{tabular}{rrrr}
\toprule
\rowcolor{navy!10}\textbf{Row}&\textbf{Ridge}&\textbf{LASSO}&\textbf{Ridge $-$ LASSO}\\
\midrule
1&-139.711&-130.368&-9.344\\
2&-112.954&-112.470&-0.485\\
3&-219.990&-218.053&-1.937\\
4&-149.057&-151.148& 2.091\\
5&-228.157&-229.846& 1.689\\
6&-102.596&-100.072&-2.524\\
\bottomrule
\end{tabular}
\end{center}

The two models give similar predictions for five rows; row 1 differs by about
$9.34$ mm. The prediction file does not contain observed
\texttt{hipcenter} values, so R can compare the predictions but cannot determine
which model is more accurate on these six cases. That requires observed test
responses or repeated cross-validation performance.

\section{Exam-ready R workflow}

\begin{enumerate}
  \item Read the files and check dimensions, names, structure, and missing values.
  \item Fit \texttt{lm(hipcenter \textasciitilde\ ., data = seatpos)}.
  \item Report coefficient interpretations, $R^2$, adjusted $R^2$, and the
  overall $F$ test.
  \item Run the four standard R diagnostic plots and calculate influence and VIF.
  \item Compare the null and full models with \texttt{anova(null\_ols, full\_ols)}.
  \item Build \texttt{X} and \texttt{X\_new} with \texttt{model.matrix}.
  \item Fix the seed and fold IDs, then use the same folds for ridge and LASSO.
  \item State explicitly whether \texttt{lambda.min} or \texttt{lambda.1se} is
  used.
  \item Extract nonzero LASSO slopes from \texttt{coef()}.
  \item Predict with both models and compare the resulting values without
  claiming a winner when test outcomes are unavailable.
\end{enumerate}

\end{document}

